\notechapter{N-20240323}{Combinatorics: Identities, Generating Functions, Catalan Numbers, and More}


\section{Binomial identities}

The note opens with several binomial identities and, more importantly, their combinatorial meanings.  The basic identity
\[
  \sum_{i=0}^n \binom ni=2^n
\]
counts all subsets of an $n$-element set: choosing a subset means choosing its size $i$ and then choosing $i$ elements.

Pascal's identity
\[
  \binom ni+\binom n{i+1}=\binom{n+1}{i+1}
\]
can be proved by fixing a distinguished element in an $(n+1)$-element set.  An $(i+1)$-subset either contains the distinguished element or it does not.  The two cases give the two terms on the left.

Vandermonde's identity
\[
  \sum_i \binom mi \binom n{r-i}=\binom{m+n}{r}
\]
counts choosing $r$ balls from two boxes with $m$ and $n$ balls.  If $i$ balls are chosen from the first box, then $r-i$ are chosen from the second.

Another useful identity is
\[
  \binom nr \binom{n-r}{k-r}=\binom nk \binom kr.
\]
Both sides count a pair $(R,K)$ with $R\subseteq K\subseteq[n]$, $|R|=r$, and $|K|=k$.

\section{Good subsets}

A nonempty subset $A\subseteq\{1,2,\ldots,n\}$ is called good if
\[
  |A|\le \min_{x\in A}x.
\]
Let $a_n$ be the number of good subsets.  The problem asks one to prove
\[
  a_{n+2}=a_{n+1}+a_n+1.
\]

A clean proof splits good subsets of $[n+2]$ into cases.

First, if $n+2\notin A$, then $A$ is just a good subset of $[n+1]$, giving $a_{n+1}$ possibilities.

Second, if $A=\{n+2\}$, this gives one extra subset.

Third, suppose $n+2\in A$ and $|A|\ge2$.  Remove $n+2$ and subtract $1$ from every remaining element.  If
\[
  A'=A\setminus\{n+2\},
\]
then the condition becomes
\[
  |A'|+1\le \min A'.
\]
After shifting $B=\{x-1:x\in A'\}$, this is exactly
\[
  |B|\le \min B,
\]
so $B$ is a good subset of $[n]$.  This gives $a_n$ possibilities.  Hence
\[
  a_{n+2}=a_{n+1}+a_n+1.
\]

The handwritten page also gives a summation proof by fixing $|A|=k$.  If $|A|=k$, then every element must be at least $k$, so one chooses the remaining data from a tail interval.  The recurrence then follows by repeated Pascal identities.

\section{Generating functions}

For a sequence $(a_n)_{n\ge0}$, its ordinary generating function is
\[
  f(x)=\sum_{n\ge0}a_nx^n.
\]
The note emphasizes a key principle: coefficients can encode counting problems.

For example, the number of nonnegative integer solutions of
\[
  x_1+x_2+\cdots+x_k=n
\]
is the coefficient of $x^n$ in
\[
  (1+x+x^2+\cdots)^k=\frac1{(1-x)^k}.
\]
Using the binomial expansion,
\[
  \frac1{(1-x)^k}=\sum_{n\ge0}\binom{n+k-1}{k-1}x^n,
\]
so the number of solutions is
\[
  \binom{n+k-1}{k-1}.
\]

The notes also use generating functions to separate parity or color constraints.  A typical trick is to multiply several series, each recording the allowed choices for one box or color, and then read off the coefficient of the target power.

% N-20240323 overflow relief A
\enlargethispage{4\baselineskip}
\section{A probability example with repeated trials}

One example considers repeated trials and asks for a probability after $n$ operations.  The handwritten solution introduces a generating function and then rewrites nested binomial sums using exponentials.  The important method is not the final closed form alone but the transformation:
\[
  \text{complicated convolution}
  \quad\longrightarrow\quad
  \text{product/composition of generating functions}.
\]
This is a standard way to turn repeated independent operations into algebra.  Binomial coefficients record choices; exponential series record independent counts; coefficient extraction recovers the desired probability.

\section{Catalan numbers}

The note then turns to Catalan numbers.  One model is the number of ways to triangulate a convex $(n+2)$-gon:
\[
  C_n=\frac1{n+1}\binom{2n}{n}.
\]
For example, a hexagon has $14$ triangulations.

Let $h_n$ denote the number of triangulations of an $(n+2)$-gon.  Choose the triangle incident with a fixed side.  If its third vertex cuts the polygon into two smaller polygons, the numbers of triangulations multiply.  Therefore
\[
  h_n=\sum_{i=0}^{n-1}h_i h_{n-1-i},\qquad h_0=1.
\]
If
\[
  H(x)=\sum_{n\ge0}h_nx^n,
\]
then the recurrence gives
\[
  H(x)=1+xH(x)^2.
\]
Solving the quadratic and choosing the branch with $H(0)=1$ gives
\[
  H(x)=\frac{1-\sqrt{1-4x}}{2x},\qquad
  h_n=C_n.
\]

This derivation explains why Catalan numbers appear in triangulations, parentheses, Dyck paths, and many noncrossing structures: all of them are recursively split into a left and a right part.

\section{Convex sets and convex hulls}

The geometry part defines a convex set $S$ by the condition that for any $x,y\in S$ and any $t\in[0,1]$,
\[
  tx+(1-t)y\in S.
\]
Equivalently, $S$ contains the line segment between any two of its points.  The note contrasts convex and non-convex examples.

For finitely many points $x_1,\ldots,x_k$, their convex hull is
\[
  \left\{\sum_{i=1}^k t_ix_i:\ t_i\ge0,\ \sum_{i=1}^k t_i=1\right\}.
\]
The handwritten proof uses induction on $k$: the case $k=2$ is the definition of a segment; the induction step writes the convex combination as a combination of one point and a convex combination of the remaining points.

\section{Separating convex sets}

The note sketches a separation idea.  For two disjoint convex sets, one tries to find a point pair realizing minimal distance and then uses the perpendicular bisector of the segment between them as a separating line.  In symbols, if $u\in S_1$ and $v\in S_2$ are closest points, then the hyperplane perpendicular to $v-u$ near the midpoint separates the two sets.

In finite-dimensional Euclidean geometry, this is the intuitive core of the separating hyperplane theorem.  The handwritten diagrams show projections onto a line and inequalities involving squared distances.  The algebra expands
\[
  \|u+t(a-u)-v\|^2
\]
for small $t$ and extracts a sign condition by letting $t\to0$.

\section{Planar graphs and Euler-type counting}

The later pages move toward graph theory and planar configurations.  The recurring idea is that a planar drawing imposes counting restrictions.  Vertices, edges, and faces cannot vary independently; they are constrained by Euler-type formulas.

This connects naturally with the previous convex-geometry material: triangulations, planar graphs, polygon decompositions, and convex hulls are different languages for controlling incidence data.

\section{Graphs and adjacency matrices}

A graph $G$ consists of a vertex set $V=\{v_1,\ldots,v_n\}$ and an edge relation.  For a simple graph, the adjacency matrix $A(G)=(a_{ij})$ is defined by
\[
  a_{ij}=
  \begin{cases}
  1,&\text{if there is an edge from }v_i\text{ to }v_j,\\
  0,&\text{otherwise.}
  \end{cases}
\]
For a multigraph, $a_{ij}$ records the number of edges from $v_i$ to $v_j$.

The source quotes the standard theorem: the $(i,j)$-entry of $A(G)^\ell$ equals the number of walks of length $\ell$ from $v_i$ to $v_j$.  The proof is matrix multiplication.  For $\ell=2$,
\[
  (A^2)_{ij}=\sum_k a_{ik}a_{kj},
\]
which counts all intermediate vertices $v_k$.  The general case follows by induction.

% N-20240323 overflow relief B
\enlargethispage{4\baselineskip}
\section{Expanded synthesis and advanced context}

This note collects several ways of counting: direct binomial identities, recurrence, generating functions, Catalan decomposition, convex geometry, and graph adjacency matrices.  The common habit is to replace a raw counting problem by a structure-preserving encoding: subsets become binomial coefficients, recursions become generating functions, triangulations become Catalan recurrences, convex hulls become linear combinations, and walks become matrix powers.



\paragraph{Expanded synthesis.}
The number-theoretic content of this chapter is best understood as arithmetic seen through algebraic structure. The earlier sections (`Binomial identities`, `Good subsets`, `Generating functions`, `A probability example with repeated trials`, `Catalan numbers`) compute divisibility, congruences, residues, prime factorizations, or finite-field examples. The deeper language is ideals, quotient rings, unit groups, valuations, and field extensions.

Divisibility in $\mathbb Z$ is containment of principal ideals; congruence modulo $n$ is equality inside the quotient ring $\mathbb Z/n\mathbb Z$; invertibility modulo $n$ means being a unit; the Chinese remainder theorem is a product decomposition of quotient rings. Prime factorization supplies coordinates through valuations:
\[
  n=\prod_p p^{v_p(n)}.
\]
Gcd and lcm are coordinatewise minimum and maximum in these valuation coordinates.

This algebraic reading explains why elementary tricks scale upward. Euclidean algorithms become the theory of Euclidean domains and PIDs; divisor sums become Dirichlet convolution and Euler products; roots of unity become cyclotomic fields and Galois actions; finite polynomial arithmetic becomes finite-field theory. The original computations should therefore be kept as the algorithmic core, while the advanced language identifies the structure that makes the algorithms work.


\section{Elementary-to-advanced bridge: Catalan numbers as recursive geometry}

Catalan numbers count many objects: correctly matched parentheses, triangulations of polygons, rooted binary trees, Dyck paths, and noncrossing partitions.  A standard formula is
\[
  C_n=\frac1{n+1}\binom{2n}{n}.
\]
The recursion
\[
  C_{n+1}=\sum_{i=0}^n C_iC_{n-i}
\]
comes from choosing the first return of a Dyck path or the root split of a binary tree.

The generating function $C(x)=\sum_{n\ge0}C_nx^n$ satisfies
\[
  C(x)=1+xC(x)^2.
\]
Solving gives
\[
  C(x)=\frac{1-\sqrt{1-4x}}{2x}.
\]
Thus an elementary recursion becomes an algebraic generating function.

\section{Elementary-to-advanced bridge: lattice paths and reflection principle}

Many Catalan formulas are proved by the reflection principle.  The number of paths from $(0,0)$ to $(n,n)$ that never go above the diagonal is
\[
  \binom{2n}{n}-\binom{2n}{n+1}=\frac1{n+1}\binom{2n}{n}.
\]
The subtracted term counts bad paths by reflecting the part up to the first diagonal crossing.

This is a model example of an involutive or bijective proof.  Rather than evaluating a sum, one constructs a symmetry pairing bad objects with easier objects.

\section{Elementary-to-advanced bridge: graph walks and adjacency matrices}

If $G$ is a finite graph with adjacency matrix $A$, then
\[
  (A^k)_{ij}
\]
counts walks of length $k$ from vertex $i$ to vertex $j$.  This turns walk counting into matrix powers.

If $A$ is diagonalizable with eigenvalues $\lambda_1,\ldots,\lambda_n$, the growth of walks is controlled by the largest eigenvalues.  In a regular connected graph, the largest eigenvalue is the degree, and the second-largest eigenvalue controls expansion and mixing.

Thus an elementary graph-walk count leads directly to spectral graph theory.

\section{Elementary-to-advanced bridge: random walks and Markov chains}

Normalize the adjacency matrix of a $d$-regular graph by
\[
  P=\frac1d A.
\]
Then $P$ is the transition matrix of a random walk.  The distribution after $k$ steps is
\[
  \pi_k=\pi_0P^k.
\]
Eigenvalues of $P$ determine convergence to stationarity.  If the graph is connected and non-bipartite, the walk converges to the uniform distribution.

The same object therefore has two readings: $A^k$ counts deterministic walks, while $P^k$ describes probabilities.  The note's graph-walk calculations are the finite beginning of Markov-chain theory.

% Second-pass deepening for waves 2-4: wave 3 A-C part 2
\section{Second-pass deepening: Catalan recursion and moduli of parenthesization}

Polygon triangulations and parenthesizations are the same combinatorial structure.  A triangulation of an $(n+2)$-gon corresponds to a way of multiplying $n+1$ factors with binary parentheses.  Choosing the triangle incident to a fixed edge splits the polygon into two smaller polygons, producing the Catalan recursion.

This is the combinatorics behind the associahedron, a polytope whose vertices are parenthesizations and whose edges are single reassociations:
\[
  (ab)c\leftrightarrow a(bc).
\]
The elementary Catalan count therefore touches the geometry of moduli spaces of bracketings.

\section{Second-pass deepening: adjacency spectra and expansion}

The theorem that $(A^k)_{ij}$ counts walks becomes much stronger when combined with eigenvalues.  If a graph is $d$-regular with adjacency eigenvalues
\[
  d=\lambda_1\ge |\lambda_2|\ge\cdots,
\]
then $|\lambda_2|$ controls how quickly random walks mix and how well the graph expands.

Expander graphs are sparse graphs that behave like highly connected graphs.  They are central in theoretical computer science, coding theory, and combinatorics.  The elementary walk-counting matrix is the first tool for studying them.
